\documentclass[11pt,a4paper]{article}

\usepackage[utf8]{inputenc}
\usepackage[T1]{fontenc}
\usepackage[spanish,es-tabla]{babel}
\usepackage{lmodern}
\usepackage[margin=2.2cm]{geometry}
\usepackage{graphicx}
\usepackage{tabularx}
\usepackage{booktabs}
\usepackage{array}
\usepackage{enumitem}
\usepackage[table]{xcolor}
\usepackage{titlesec}
\usepackage{setspace}
\usepackage[hidelinks]{hyperref}

\definecolor{espeverde}{RGB}{0,102,51}
\titleformat{\section}{\large\bfseries\color{espeverde}}{\thesection.}{0.6em}{\MakeUppercase}
\titleformat{\subsection}{\normalsize\bfseries}{\thesubsection}{0.6em}{}
\setlength{\parskip}{0.5em}
\setlength{\parindent}{0pt}
\setstretch{1.15}
\setlist{itemsep=0.15em, topsep=0.3em}

\begin{document}

% ============================================================
% PORTADA
% ============================================================
\begin{titlepage}
    \centering
    \includegraphics[width=0.28\textwidth]{img/LogoESPE.png}\par\vspace{0.8cm}
    {\Large\bfseries UNIVERSIDAD DE LAS FUERZAS ARMADAS -- ESPE\par}
    \vspace{0.4cm}
    {\large Departamento de Ciencias de la Computación\par}
    {\large Carrera de Ingeniería de Software\par}
    \vspace{1cm}
    {\large\bfseries Desarrollo de Software Global Aplicado a la Interculturalidad\par}
    \vspace{1cm}
    {\large FORMATO 1 -- PROPUESTA DEL TEMA DEL PROYECTO\par}
    \vspace{0.5cm}
    {\LARGE\bfseries\color{espeverde} Atlas Anatómico Inmersivo:\\[0.2cm]
    laboratorio de Realidad Virtual bilingüe (español--kichwa) para el estudio del ojo, los pulmones y el riñón\par}
    \vspace{1cm}
    \begin{tabular}{rl}
        \textbf{Grupo:} & 8 \\[0.2cm]
        \textbf{Integrantes:} & Germán Cáceres \\
                              & Carlos Rivera \\
                              & Samir Mideros \\[0.2cm]
        \textbf{Docente:} & [Nombre del docente] \\[0.2cm]
        \textbf{Fecha:} & 25 de octubre de 2026 \\
    \end{tabular}
    \vfill
    {\large Sangolquí -- Ecuador\par}
\end{titlepage}

% ============================================================
\section{Título del proyecto}
% ============================================================
\textbf{Atlas Anatómico Inmersivo: laboratorio de Realidad Virtual bilingüe (español--kichwa) para el estudio del ojo, los pulmones y el riñón.}

% ============================================================
\section{Introducción}
% ============================================================
La enseñanza de la anatomía humana se apoya tradicionalmente en láminas, libros y, en el mejor de los casos, en maquetas físicas o en el trabajo con cadáveres. Estos recursos son costosos, escasos fuera de las grandes ciudades y difíciles de usar para comprender la estructura tridimensional de los órganos: un estudiante que observa un dibujo plano del ojo o del riñón debe \emph{imaginar} cómo se superponen sus capas internas.

Las tecnologías inmersivas, en particular la Realidad Virtual (RV), permiten que el estudiante se ubique \emph{dentro} de un laboratorio, tome un órgano con sus manos, lo gire, lo amplíe y lo separe en sus partes. Diversos estudios reportan que la RV mejora la comprensión espacial y la motivación en la enseñanza de la anatomía \cite{moro2017}.

Sin embargo, la mayoría de atlas anatómicos digitales están disponibles solo en inglés o en español y no consideran la diversidad lingüística del Ecuador, un Estado constitucionalmente intercultural y plurinacional donde el kichwa es idioma oficial de relación intercultural \cite{constitucion}. El presente proyecto propone un \textbf{Atlas Anatómico Inmersivo} desarrollado con \textbf{Blender} y \textbf{Unity} que presenta los contenidos en español y kichwa, con narración por voz, subtítulos y modos de interacción accesibles. De esta forma, el proyecto une las tecnologías inmersivas con el software inclusivo e intercultural, ejes de la asignatura.

% ============================================================
\section{Planteamiento del problema}
% ============================================================
\textbf{Situación actual.} En el Ecuador, los estudiantes de bachillerato en Ciencias y de los primeros niveles de carreras de la salud aprenden anatomía principalmente con material impreso. Las instituciones rurales e interculturales bilingües rara vez cuentan con maquetas, laboratorios o software especializado, y los recursos digitales existentes son monolingües y no son accesibles para personas con discapacidad.

\textbf{Problema existente.} No existe un recurso didáctico tridimensional, inmersivo, accesible y culturalmente pertinente que permita estudiar los órganos del cuerpo humano en español y kichwa.

\textbf{Consecuencias.} Si el problema persiste:
\begin{itemize}
    \item Los estudiantes memorizan nombres sin comprender la estructura espacial de los órganos.
    \item Se amplía la brecha educativa entre zonas urbanas y rurales.
    \item Los estudiantes kichwa hablantes aprenden ciencia en un idioma que no es el propio, lo que debilita su lengua y su identidad.
    \item Las personas con discapacidad visual, auditiva o motora quedan excluidas de las herramientas tecnológicas de aprendizaje.
\end{itemize}

\textbf{Beneficiarios afectados.} Estudiantes de bachillerato y de los primeros semestres de carreras de la salud, especialmente de instituciones del Sistema de Educación Intercultural Bilingüe, y sus docentes.

% ============================================================
\section{Justificación}
% ============================================================
\subsection*{Justificación social}
El proyecto beneficia a estudiantes y docentes que no tienen acceso a laboratorios de anatomía, y en particular a la población kichwa hablante, al ofrecer contenido científico en su lengua. Además, incorpora criterios de accesibilidad para que estudiantes con distintas capacidades puedan usarlo.

\subsection*{Justificación tecnológica}
Se utilizará \textbf{Blender} para el modelado 3D de los órganos y \textbf{Unity} con \textbf{XR Interaction Toolkit} y \textbf{OpenXR} para construir el laboratorio virtual, compatible con visores autónomos (Meta Quest 2/3) y con un modo de escritorio para quienes no disponen de visor. Estas herramientas son gratuitas para uso educativo, lo que favorece su réplica en otras instituciones.

\subsection*{Justificación académica}
El proyecto aplica los contenidos de la asignatura: desarrollo de software para contextos culturales diversos, localización e internacionalización, diseño centrado en el usuario (ISO 9241-210) \cite{iso9241}, pautas de accesibilidad (WCAG 2.2) \cite{wcag22} y trabajo colaborativo con control de versiones (Git).

% ============================================================
\section{Objetivo general}
% ============================================================
Desarrollar un prototipo de laboratorio de Realidad Virtual con Blender y Unity que permita visualizar, manipular y estudiar modelos 3D del ojo, los pulmones y el riñón en español y kichwa, aplicando criterios de accesibilidad e inclusión.

% ============================================================
\section{Objetivos específicos}
% ============================================================
\begin{enumerate}
    \item Identificar las necesidades de aprendizaje y de accesibilidad de la población objetivo mediante una encuesta y una revisión bibliográfica.
    \item Modelar en Blender tres órganos (ojo, pulmones y riñón) separados en sus estructuras principales, optimizados para RV (menos de 50\,000 polígonos por órgano).
    \item Implementar en Unity un laboratorio virtual con interacciones de agarre, rotación, escala, vista explosionada y etiquetas bilingües (español--kichwa).
    \item Incorporar opciones de accesibilidad: narración por voz, subtítulos, alto contraste, tamaño de texto ajustable y modo sentado.
    \item Evaluar la usabilidad del prototipo con al menos cinco usuarios mediante el cuestionario SUS (\emph{System Usability Scale}).
\end{enumerate}

% ============================================================
\section{Población objetivo}
% ============================================================
\begin{itemize}
    \item \textbf{Quiénes:} estudiantes de bachillerato en Ciencias y de los primeros semestres de carreras de la salud (Medicina, Enfermería, Biología); como usuarios secundarios, sus docentes.
    \item \textbf{Edad:} entre 15 y 22 años.
    \item \textbf{Contexto:} instituciones educativas del Ecuador, incluidas las del Sistema de Educación Intercultural Bilingüe, con acceso limitado a laboratorios de anatomía.
    \item \textbf{Necesidades específicas:} contenido en su lengua materna (español o kichwa), explicaciones claras y breves, interacción sencilla para usuarios sin experiencia en RV y alternativas para personas con discapacidad visual, auditiva o motora.
\end{itemize}

% ============================================================
\section{Propuesta de solución}
% ============================================================
\subsection{¿Qué se desarrollará?}
Un laboratorio virtual con una mesa de estudio sobre la que el usuario elige un órgano. Al seleccionarlo, el órgano aparece a escala ampliada y puede tomarse con las manos (controles), girarse, ampliarse y \emph{explosionarse} para ver sus partes internas. Cada parte tiene una etiqueta en español y kichwa y una ficha con una descripción breve narrada por voz.

\subsection{Alcance del prototipo}
Para asegurar un producto funcional en el semestre, el alcance se limita a \textbf{tres órganos}, uno por integrante. Se eligieron por su geometría sencilla (pueden construirse a partir de esferas, cápsulas y curvas), por pertenecer a sistemas distintos y por su valor didáctico.

\begin{table}[h]
\centering
\small
\rowcolors{2}{gray!10}{white}
\begin{tabularx}{\textwidth}{>{\raggedright}p{2.4cm} >{\raggedright}p{2.2cm} >{\raggedright\arraybackslash}X >{\raggedright\arraybackslash}p{3.4cm}}
\toprule
\textbf{Responsable} & \textbf{Órgano (sistema)} & \textbf{Estructuras a modelar} & \textbf{Interacción especial} \\
\midrule
Germán Cáceres & Ojo -- \emph{ñawi} (sensorial) & Esclerótica, córnea, iris, pupila, cristalino, humor vítreo, retina y nervio óptico & Vista explosionada por capas \\
Carlos Rivera & Pulmones -- \emph{surku} (respiratorio) & Pulmón derecho (3 lóbulos), pulmón izquierdo (2 lóbulos), tráquea, bronquios principales y diafragma & Animación de respiración \\
Samir Mideros & Riñón -- \emph{rurun} (urinario) & Corteza, médula (pirámides renales), pelvis renal, arteria y vena renal, uréter & Corte transversal interactivo \\
\bottomrule
\end{tabularx}
\caption{Alcance del prototipo: órganos y responsables.}
\end{table}

Cada integrante modela su órgano en Blender y desarrolla su estación en Unity; el laboratorio común (escena, menús, localización y accesibilidad) se construye de forma conjunta.

\textbf{Fuera del alcance de esta versión:} corazón, cerebro, sistema esquelético y muscular completos, simulaciones fisiológicas complejas, modo multijugador y evaluación clínica. Estos elementos se consideran trabajo futuro.

\subsection{¿Cómo funcionará?}
\begin{enumerate}
    \item El usuario elige el idioma (español o kichwa) y las opciones de accesibilidad en el menú inicial.
    \item En el laboratorio, selecciona un órgano en la mesa de estudio.
    \item Manipula el órgano: agarrar, rotar, escalar y activar la vista explosionada, el corte o la animación.
    \item Al señalar una estructura, se resalta y se muestra su nombre bilingüe; puede escuchar su descripción.
    \item Opcionalmente, resuelve un cuestionario corto de identificación de estructuras.
\end{enumerate}

\subsection{Tecnologías}
\begin{itemize}
    \item \textbf{Blender 4.x:} modelado, texturizado y animación de los órganos; exportación en formato FBX/glTF.
    \item \textbf{Unity 6 (LTS):} motor del laboratorio, con \textbf{XR Interaction Toolkit} y \textbf{OpenXR}.
    \item \textbf{Meta Quest 2/3} como visor objetivo y \textbf{XR Device Simulator} para el modo de escritorio.
    \item \textbf{Unity Localization} para la gestión de textos en español y kichwa.
    \item \textbf{Git y GitHub} para el control de versiones y el trabajo colaborativo.
    \item \textbf{Z-Anatomy / BodyParts3D} (licencias abiertas) como referencia anatómica \cite{bodyparts3d}.
\end{itemize}

% ============================================================
\section{Componente de interculturalidad}
% ============================================================
\begin{itemize}
    \item \textbf{Grupo cultural beneficiado:} estudiantes de los pueblos y nacionalidades kichwa del Ecuador, sin dejar de lado a los estudiantes hispanohablantes.
    \item \textbf{Respeto a la diversidad cultural:} todos los textos, etiquetas y narraciones estarán disponibles en español y kichwa, con el mismo nivel de detalle en ambos idiomas; el kichwa no será una traducción secundaria ni incompleta.
    \item \textbf{Prevención de sesgos y estereotipos:} no se usarán vestimentas, rostros ni símbolos culturales como decoración. Los términos en kichwa (por ejemplo \emph{ñawi}, \emph{surku}, \emph{rurun}) se validarán con un hablante nativo o docente de Educación Intercultural Bilingüe y con el diccionario oficial del Ministerio de Educación \cite{kichwa}.
    \item \textbf{Representación adecuada de la información:} el contenido anatómico será científico y neutral. Cuando un término no tenga equivalente en kichwa, se indicará de forma explícita y se mantendrá el término técnico en español, sin inventar palabras.
\end{itemize}

% ============================================================
\section{Componente de accesibilidad e inclusión}
% ============================================================
\begin{description}
    \item[Accesibilidad visual:] paneles con contraste de al menos 4.5:1 \cite{wcag22}, tamaño de texto ajustable, colores de resaltado distinguibles para personas con daltonismo (apoyados por contornos, no solo por color) y narración por voz de todas las fichas.
    \item[Accesibilidad auditiva:] subtítulos sincronizados de toda narración, indicadores visuales y vibración de los controles en lugar de señales solo sonoras.
    \item[Accesibilidad cognitiva:] lenguaje claro y frases cortas, un solo órgano a la vez, tutorial guiado al inicio y navegación de máximo dos niveles de menú.
    \item[Accesibilidad motora:] modo sentado, ajuste de la altura de la mesa, interacción a distancia por rayo (sin necesidad de alcanzar objetos), botones de gran tamaño, teletransporte en lugar de desplazamiento continuo para reducir el mareo, y modo de escritorio con teclado y ratón para usuarios sin visor.
\end{description}

% ============================================================
\section{Resultados esperados}
% ============================================================
\begin{itemize}
    \item Tres modelos 3D (ojo, pulmones y riñón) optimizados para RV, con sus estructuras separadas.
    \item Prototipo funcional del laboratorio en Unity, ejecutable en Meta Quest y en modo de escritorio.
    \item Glosario bilingüe español--kichwa de las estructuras anatómicas, validado.
    \item Manual de usuario.
    \item Informe de evaluación de usabilidad (SUS).
    \item Informe final del proyecto y repositorio en GitHub.
\end{itemize}

% ============================================================
\section{Cronograma inicial}
% ============================================================
\begin{table}[h]
\centering
\rowcolors{2}{gray!10}{white}
\begin{tabular}{l c c c c}
\toprule
\textbf{Actividad} & \textbf{Mes 1} & \textbf{Mes 2} & \textbf{Mes 3} & \textbf{Mes 4} \\
\midrule
Investigación y definición del alcance & X & & & \\
Diseño (bocetos, glosario bilingüe, interfaz) & X & X & & \\
Modelado de órganos en Blender & & X & & \\
Desarrollo del laboratorio en Unity & & X & X & \\
Localización y accesibilidad & & & X & \\
Pruebas de usabilidad & & & X & X \\
Presentación final & & & & X \\
\bottomrule
\end{tabular}
\caption{Cronograma inicial del proyecto.}
\end{table}

% ============================================================
\section{Bibliografía inicial}
% ============================================================
\begingroup
\renewcommand{\section}[2]{}
\begin{thebibliography}{9}
\bibitem{moro2017} Moro, C., Štromberga, Z., Raikos, A. y Stirling, A. (2017). The effectiveness of virtual and augmented reality in health sciences and medical anatomy. \emph{Anatomical Sciences Education}, 10(6), 549--559.
\bibitem{bodyparts3d} Mitsuhashi, N., Fujieda, K., Tamura, T., Kawamoto, S., Takagi, T. y Okubo, K. (2009). BodyParts3D: 3D structure database for anatomical concepts. \emph{Nucleic Acids Research}, 37, D782--D785.
\bibitem{wcag22} World Wide Web Consortium (W3C). (2023). \emph{Web Content Accessibility Guidelines (WCAG) 2.2}. \url{https://www.w3.org/TR/WCAG22/}
\bibitem{iso9241} International Organization for Standardization. (2019). \emph{ISO 9241-210:2019. Ergonomics of human-system interaction -- Part 210: Human-centred design for interactive systems}.
\bibitem{unityxr} Unity Technologies. (2026). \emph{XR Interaction Toolkit documentation}. \url{https://docs.unity3d.com/Packages/com.unity.xr.interaction.toolkit@latest}
\bibitem{blender} Blender Foundation. (2026). \emph{Blender Reference Manual}. \url{https://docs.blender.org/manual/es/latest/}
\bibitem{kichwa} Ministerio de Educación del Ecuador. (2009). \emph{Kichwa Yachakukkunapak Shimiyuk Kamu: diccionario kichwa--castellano}. Quito: DINEIB.
\bibitem{constitucion} Asamblea Constituyente. (2008). \emph{Constitución de la República del Ecuador}, art. 1 y art. 2.
\end{thebibliography}
\endgroup

\end{document}
